% Research contribution fragment for the parent article.
% Source baseline: ProveIt da2792fc9cc8ef19c6a314be6c65e1437fbdc66b.
% All labels use mt10:. This is not a standalone final article.

\section{Completion of the smaller-order-ten reconstruction theorem}
\label{mt10:section}

Two incoming reports leave complementary parts of the same problem open.
The \emph{Alternating Tails and Spectral Defects} report \cite{RepoAlternating} proves the
smaller-order-ten statement outside one residue class modulo $2048$;
the \emph{Quartic Spectral Reconstruction} report \cite{RepoQuartic} proves it for all
sufficiently large larger orders. Their thresholds leave exactly one
matrix unaccounted for. We prove that this matrix is nonsingular by a
short exact certificate and hence complete the entire smaller-order-ten
family. The finite polynomial identities and the sufficient threshold
used from those reports are replayed independently below. Their
underlying reductions remain attributed to those reports.

This completion concerns formal rational coordinate rank. It implies
identities among actual polygamma sums, but does not assert arithmetic
independence of evaluated zeta values. The unrestricted unequal-order
conjecture remains open.

\subsection{Moment coordinates and the precise theorem}

For integers $s\ge2$ put
\[
 A_s(n)=\zeta(s,n+1)=\zeta(s)-H_n^{(s)},\qquad x_n=n+\tfrac12.
\]
For an integer $b\ge6$, initially define
\[
 D_{10,2b}(w)=\sum_{n\ge0}A_{10}(n)A_{2b}(n)x_n^{-w},
 \qquad \operatorname{Re}w>-2b-7.
\]
The centered Hurwitz expansion contains only powers
$x_n^{-2b-8-2j}$, so continuation by subtraction is regular at every
nonpositive even integer. Write $D_r=D_{10,2b}(-2r)$. Define the
mutually inverse finite transforms
\begin{align}
 C_j&=\sum_{r=0}^j\binom{2j+1}{2r}2^{-2(j-r)}D_r,
 \label{mt10:C}\\
 D_r&=\sum_{j=0}^r\frac{\binom{2r}{2j}}{2j+1}
                 B_{2r-2j}(1/2)C_j.\label{mt10:D}
\end{align}
Finite summation by parts gives, for every $j\ge b$,
\[
 C_j=\sum_{k=0}^{b-1}(M^{\infty}_{5,b})_{j-b,k}Z_k,
 \qquad Z_k=\zeta(2k),\quad Z_0=-\tfrac12,
\]
where, for $t\ge0$ and $0\le k<b$,
\begin{equation}\label{mt10:M}
 (M^{\infty}_{5,b})_{tk}
 =\mu_{t+k-4}\left\{
 \frac{(2b+2t-9)!}{(2b-2k)!}
 +\mathbf1_{k<5}\frac{(2t+1)!}{(10-2k)!}\right\},
 \qquad
 \mu_j=\begin{cases}B_{2j}/(2j)!,&j\ge0,\\0,&j<0.\end{cases}
\end{equation}
Let $M_{5,b}$ be its first $b$ rows. Equivalently, each entry is
\[
 B_{2(t+k-4)}\left\{
 \frac{\binom{2t+2}{10-2k}}{2t+2}
 +\frac{\binom{2b+2t-8}{2b-2k}}{2b+2t-8}\right\},
\]
where a negative Bernoulli index gives zero, and a binomial summand
is omitted unless its lower index lies between two and its upper
index. This last convention is essential.

The formal coordinate space here has independent labels
$C_0,\ldots,C_{b-1},Z_0,\ldots,Z_{b-1}$. Its evaluation map assigns
the indicated complex numbers to those labels. In these coordinates
the first $2b$ transformed moments have matrix
$\operatorname{diag}(I_b,M_{5,b})$. Since the diagonal of
\eqref{mt10:D} is $1/(2r+1)$, the first $2b$ original moments have
determinant $\det M_{5,b}/(4b-1)!!$.

\begin{theorem}[Every larger even order at smaller order ten]
\label{mt10:main}
For every integer $b\ge6$, $M_{5,b}$ is nonsingular. Consequently
\[
 D_{10,2b}(0),D_{10,2b}(-2),\ldots,D_{10,2b}(-2(2b-1))
\]
have full formal coordinate rank $2b$, and every subsequent centered
even moment is an explicit rational linear combination of this prefix.
\end{theorem}

\subsection{The five-dimensional contraction and its normalization}

For $0\le i,\ell\le4$ set
\[
 n_i=b-5+i,\qquad N_{i\ell}=2(n_i-\ell)+1,
 \qquad R_\ell=2b-8+2\ell,
\]
\[
 F_{i\ell}=\frac1{(2\ell+1)!}
       \prod_{r=0}^{2\ell-1}(2n_i-r)(2n_i+3+r).
\]
An empty product is one. Define finite integer-polynomial coefficients
$d_{ih}$ by the Newton expansion
\begin{equation}\label{mt10:newton}
 (2n_i+2-y)(4n_i+2-y)^{\underline{2i}}
   =\sum_{h=0}^{2i+1}d_{ih}y^{\underline h}.
\end{equation}
Thus explicitly
\[
 d_{ih}=\frac1{h!}\sum_{j=0}^h(-1)^{h-j}\binom hj
           (2n_i+2-j)(4n_i+2-j)^{\underline{2i}}.
\]
For independent variables $U_1,\ldots,U_9$, let
\begin{equation}\label{mt10:A}
 \mathcal A_{i\ell}=2F_{i\ell}+
 \frac{R_\ell!}{(2n_i+1)(2n_i+2)(2\ell+2)!N_{i\ell}!}
 \sum_{k=0}^i d_{i,2k+1}N_{i\ell}^{\underline{2k+1}}
                      U_{5-i+\ell+k}.
\end{equation}
Ratios of factorials have fixed index difference and mean the
corresponding finite rational products in $b$.

For $b\ge10$, the exact reduction is
\begin{equation}\label{mt10:ratio}
 \rho_{5,b}:=\frac{\det M_{5,b}}{\det R\det E\det H_{5,b}}
  =\frac{\det\mathcal A(b;2B_{2b-2},\ldots,2B_{2b-18})}
         {\det F(b)},
\end{equation}
where
\[
 R_{tt}=(2b+2t-9)!,\qquad E_{kk}=\frac1{(2b-2k)!},
 \qquad (H_{5,b})_{tk}=\mu_{t+k-4}.
\]
All three determinants in the denominator are nonzero.

Here are details sufficient to check the contraction without an
assumed inverse-Hankel formula. The monic polynomial rows
\[
 \pi_n(x)=\sum_{t=0}^n
 \frac{2(2n+2t+1)!(2n+1)!}
 {(4n+2)!(2n-2t+1)!(2t)!}x^t
\]
are orthogonal for the functional $\mathcal N(x^r)=\mu_{r+1}$:
\[
 \mathcal N(x^j\pi_n)=0\quad(0\le j<n),\qquad
 \mathcal N(x^n\pi_n)=h_n,
 \quad h_n=\frac{((2n+1)!)^2}{(4n+2)!(4n+3)!}>0.
\]
These follow by coefficient comparison in the elementary Pad\'e identity
\[
 e^tP_N(-t)-P_N(t)
 =\frac{(-1)^N(N!)^2}{(2N)!(2N+1)!}t^{2N+1}
        +O(t^{2N+2}),
\]
\[
 P_N(t)=\sum_{k=0}^N
       \frac{(2N-k)!N!}{(2N)!k!(N-k)!}t^k,
 \qquad N=2n+1.
\]
The coefficient cancellation through degree $2N$ is the finite binomial
identity $\sum_k(-1)^k\binom Nk(1+x)^{2N-k}=x^N(1+x)^N$;
the next coefficient is an elementary beta integral.

Replacing every row of $H_{5,b}+U^{\mathrm{pert}}$ by these monic
rows is a unit lower triangular operation. Its final five rows vanish
in columns $k\ge5$. The unperturbed retained entries, after column
reversal and division by
$c_n=(2n+1)(2n+2)!/(4n+2)!$, are precisely $F_{i\ell}$.
For the perturbation apply
\[
 \sum_{j=0}^N\binom Nj B_j(N-j)^{\underline h}
                  =N^{\underline h}B_{N-h}(1)
\]
to \eqref{mt10:newton}. Because $b\ge10$, every surviving index
$N-h$ is at least two. The even $h$ give vanishing odd Bernoulli
numbers; restoring the omitted term $j=1$, with $B_1=-1/2$, contributes
a second $F_{i\ell}$. This proves \eqref{mt10:A}, including both
factors of two. Carrying the same row operation out on $H_{5,b}$
alone gives \eqref{mt10:ratio}.

For completeness, $H_{5,b}$ is nonsingular: after conjugating by
$\operatorname{diag}(1,-1,1,\ldots)$ and splitting off its first
three indices, it has the form
\[
 \begin{pmatrix}ee^{\mathsf T}&B\\B^{\mathsf T}&-P\end{pmatrix},
 \qquad e=(0,0,1)^{\mathsf T},
\]
where $P$ is positive definite. Its entries are moments on the infinite
positive support $\{(2\pi m)^{-2}:m\ge1\}$. The first two rows
of $B$ have a nonsingular antitriangular two-by-two minor. Therefore
$ee^{\mathsf T}+BP^{-1}B^{\mathsf T}$ is positive definite, which
proves nonsingularity by the Schur complement. This argument also
shows that the unperturbed contraction $F$ is nonsingular.

\subsection{The inherited congruence, with its finite certificate replayed}

Define
\begin{align*}
 D_{10}(b)={}&4096b^{12}-110592b^{11}+1236992b^{10}-7303680b^9\\
 &+23705088b^8-38140416b^7+14925056b^6+24490080b^5\\
 &-4469984b^4-15838992b^3-2571948b^2+4244400b+1488375.
\end{align*}
The exact five-by-five polynomial identities are
\begin{align}
 \det F={}&\frac{16384}{4465125}
 (b-3)(b-2)^2(b-1)(2b-7)(2b-5)^2(2b-3)\notag\\
 &\qquad\cdot(4b-15)(4b-13)(4b-11)^2(4b-9)^2(4b-7)^2
 (4b-5)(4b-3)D_{10}(b),\label{mt10:Fdet}\\
 \frac{\det\mathcal A}{\det F}
   &=\frac{P_{10}(b;U_1,\ldots,U_9)}{
        138240(2b-7)(2b-5)D_{10}(b)},\label{mt10:P}
\end{align}
where $P_{10}$ is an integer polynomial. Equations \eqref{mt10:A}
and \eqref{mt10:P} define it exactly; its supplied sparse certificate
contains $10{,}403$ monomials, grouped into $439$ monomials in the
$U$ variables. Their coefficients have degree at most $24$ in $b$.

The replay forms the row denominators of \eqref{mt10:A}, constructs
the $120$ signed permutation products of the cleared matrix, and
cross-multiplies \eqref{mt10:P} in
$\mathbb Q[b,U_1,\ldots,U_9]$. Equality is checked coefficient by
coefficient. Thus polynomial integrality and the displayed determinant
identity are finite rational certificates, without interpolation.

Put
\[
 \mathcal B(x)=1+20068x-26384x^2+2816x^3-19712x^4
                         +4096x^5-24576x^6.
\]
The uniform Bernoulli congruence is
\begin{equation}\label{mt10:B16}
           2B_{2m}\equiv\mathcal B(m)\pmod {65536}
                  \quad(m\ge12).
\end{equation}
All congruences are in $\mathbb Z_{(2)}$, the rationals with odd
denominator. One direct proof starts with Faulhaber's formula at
$256$. The terms other than $256B_{2m}$ have valuation at least
$23$; after division by $128$, the even summands vanish at the
required precision because $2m\ge24$. For odd $u$, $8\mid u^2-1$,
so expansion of $u^{2m}=(1+(u^2-1))^m$ gives
\[
 2B_{2m}\equiv\frac1{128}
       \sum_{u=1,3,\ldots,255}\sum_{j=0}^7
                    \binom mj(u^2-1)^j\pmod {65536}.
\]
The omitted terms have valuation at least $3\cdot8-7=17$.
Expanding this finite average in $\mathbb Q[m]$ and reducing every
coefficient modulo $65536$ gives \eqref{mt10:B16}. The replay
checks this coefficientwise, including the denominators of
$\binom mj$. The bound $v_2(B_{2j})\ge-1$ used in the Faulhaber
step follows inductively from the odd-denominator recurrence
\eqref{mt10:Bernoulli-recurrence} below.

The finite polynomial substitution then gives
\begin{equation}\label{mt10:congruence}
 P_{10}\bigl(b;\mathcal B(b-1),\ldots,\mathcal B(b-9)\bigr)
                  \equiv32\mathcal R(b)\pmod {65536},
\end{equation}
where
\[
 \mathcal R(b)=-512b^8+128b^7-928b^6-816b^5+464b^4
                         +352b^3+118b^2+217b+694.
\]
The replay verifies every coefficient of \eqref{mt10:congruence}
using modular polynomial multiplication. The linear coefficient of
$\mathcal R$ is odd and every coefficient of degree at least two
is even. Consequently
\[
 v_2(\mathcal R(x)-\mathcal R(y))=v_2(x-y)\quad(x\ne y).
\]
Indeed the quotient by $x-y$ is odd. Direct substitution gives
$\mathcal R(962)\equiv0\pmod {2048}$, so $962$ is its unique
root modulo $2048$.

Since $D_{10}(b)$ is odd at every integer $b$, the denominator in
\eqref{mt10:P} has valuation $10$. For $b\ge21$ all nine Bernoulli
indices are in the range of \eqref{mt10:B16}. If
$b\not\equiv962\pmod {2048}$, equations \eqref{mt10:B16}--
\eqref{mt10:congruence} therefore prove
\begin{equation}\label{mt10:valuation}
                 v_2(\rho_{5,b})=v_2(b-962)-5.
\end{equation}
For $6\le b\le20$, direct rational determinants of the original
matrix \eqref{mt10:M} give the respective valuations
\[
 -4,-4,-6,-5,-2,-5,-4,-5,-3,-5,-4,-5,-1,-5,-4.
\]
The direct matrices, their determinants, and their normalizations are
included in the replay record. In particular, the contraction is never
used below its valid threshold $b\ge10$.

\subsection{A transparent sufficient large-order threshold}

We specialize the incoming eventual-rank proof to $a=5$ and replace
its large rational constant by an elementary upper bound. This makes
the threshold independently reviewable.

Define
\[
 J(b)=\prod_{i=0}^4(2n_i+1)(2n_i+2),\qquad
 E(b)=\prod_{i,\ell=0}^4
             \prod_{r=1}^{\max(0,2i-4\ell-1)}(R_\ell+r),
 \quad D(b)=10!J(b)E(b).
\]
Here $\deg J=10$, $\deg E=20$, and $D\mathcal A$ has integer
polynomial entries. Put $\mathscr P=\det(D\mathcal A)$ and let
$\mathscr P_{[5]}$ be its homogeneous part of degree five in the
$U$ variables. Every entry has degree at most $46$ in $b$; hence
$\deg_b\mathscr P\le230$. The largest coefficient norm among the
$25$ entries, with all coefficients counted in
$\mathbb Z[b,U_1,\ldots,U_9]$, is
\[
 E^*=602847821074996042728613866361551228764160<10^{42}.
\]
The separate verifier expands the $25$ entries, records all $25$
norms and degrees, and reproduces this integer. The full determinant
coefficient norm is at most $K=5!(E^*)^5<10^{213}$.

Let $F_0=(2b-18)!$ and $V_j=U_jF_0/(2b-2j)!$. From
\eqref{mt10:newton},
$d_{i1}=-(i+1)(4n_i+2)^{\underline{2i}}$. Cancelling factorials
in \eqref{mt10:A} gives
\[
 \mathcal A-2F=
 \operatorname{diag}\left(\frac1{(2n_i+1)(2n_i+2)}\right)
 L(V_{5-i+\ell})_{0\le i,\ell\le4}
 \operatorname{diag}\left(\frac{R_\ell!}{F_0(2\ell+2)!}\right),
\]
where $L_{i,i-k}=d_{i,2k+1}$ is lower triangular. Therefore
\[
 \mathscr P_{[5]}=-\Gamma(b)\det(V_{i+j+1})_{0\le i,j\le4},
\]
with the explicit positive factor
\begin{align*}
 \Gamma(b)={}&\kappa J(b)^4E(b)^5
 \prod_{i=0}^4(4n_i+2)^{\underline{2i}}
 \prod_{\ell=0}^4\prod_{r=1}^{10+2\ell}(2b-18+r),\\
 \kappa={}&\frac{5!(10!)^5}{2!4!6!8!10!}
          =14932726824960000000>10^{19}.
\end{align*}
Every linear factor is at least $b$ for $b\ge20$, so
$\Gamma(b)\ge\kappa b^{230}$.

At the actual doubled Bernoulli values, Euler's even-zeta formula
and row/column reversal yield
\[
 |\mathscr P_{[5]}|
 =\frac{\Gamma(b)(4F_0)^5}{(2\pi)^{10(b-5)}}
       \det\bigl(\zeta(2b-18+2i+2j)\bigr)_{0\le i,j\le4}.
\]
The determinant on the right is a positive moment determinant on
the points $m^{-2}$. Keeping the support points $1,1/4,1/9,1/16,1/25$
in Cauchy--Binet proves
\begin{gather*}
 \det\bigl(\zeta(2b-18+2i+2j)\bigr)\ge V\,120^{-2b+18},\\
 V=\prod_{1\le r<s\le5}(r^{-2}-s^{-2})^2
        =\frac{2401}{36864000000000000}>10^{-14}.
\end{gather*}

For $b\ge30$, all nine actual variables have absolute value at most
\[
       T_b=\frac{8(2b)^{16}F_0}{(2\pi)^{2b-18}}\ge1.
\]
Thus $|\mathscr P_{[\le4]}|\le Kb^{230}T_b^4$.
After cancellation of the common powers of $2\pi$, using
$1<2\pi<8$ gives
\begin{equation}\label{mt10:dominance}
 \frac{|\mathscr P_{[\le4]}|}{|\mathscr P_{[5]}|}
 \le C\frac{b^{64}960^{2b}}{(2b-18)!},\qquad
 C=\frac{K2^{142}}{\kappa4^5V120^{18}}<10^{251}<10^{256}.
\end{equation}
The powers here can be checked directly: the uncancelled power of
$2\pi$ is $2b+22$, and $(2b)^{64}$ comes from $T_b^4$.
The bounds $K<10^{213}$, $2^{142}<10^{43}$, $\kappa>10^{19}$,
and $V>10^{-14}$ already give the stated coarse bound on $C$.

The remaining comparison is an exact integer inequality:
\begin{equation}\label{mt10:base}
       2^{384}\,10^{256}\,1920^{64}\,960^{3840}<3822!.
\end{equation}
It is verified by integer multiplication, with both sides explicitly
specified. Moreover, for $b\ge1920$,
\[
 \frac{(b+1)^{64}960^{2b+2}/(2b-16)!}
      {b^{64}960^{2b}/(2b-18)!}
 =\left(1+\frac1b\right)^{64}
            \frac{960^2}{(2b-16)(2b-17)}<\frac34.
\]
Indeed $(1+1/b)^{64}<3$ and both denominator factors are at least
$b\ge1920=2\cdot960$. It follows from
\eqref{mt10:dominance}--\eqref{mt10:base} that the leading homogeneous
term cannot be cancelled when $b\ge1920$. Since $D(b)>0$ there,
\eqref{mt10:ratio} proves nonsingularity throughout that range.

\subsection{The single missing matrix has a small exact certificate}

The congruence covers every $6\le b<1920$ except $b=962$.
At this value define
\[
 S=\operatorname{diag}(2,1,64,1/8,2)\,
 \mathcal A\,\operatorname{diag}(1,1,1/32,1/16,1/512).
\]
Thus $\det S=2^{-13}\det\mathcal A$.
Every coefficient of this affine matrix in $U_1,\ldots,U_9$ has
$2$-adic valuation at least $-3$. Therefore its reduction modulo
$2^{20}$ is determined by the inputs modulo $2^{23}$.

The nine input residues, in increasing order of the index $j$, are
\[
 (660565,5568257,602509,3180025,6422085,3501169,
                         2681469,1596009,6743605)\pmod {2^{23}}.
\]
They are obtained by the elementary Bernoulli recurrence
\begin{equation}\label{mt10:Bernoulli-recurrence}
 (2m+1)(2B_{2m})
   =2m-1-\sum_{j=1}^{m-1}\binom{2m+1}{2j}(2B_{2j}),\qquad m\ge1.
\end{equation}
Every divisor $2m+1$ is odd. The recurrence is valid modulo every
power of two and produces the residues through $m=961$ without
floating arithmetic. A separate exact rational Bernoulli evaluation
agrees with all nine inputs.

Substitution in \eqref{mt10:A} gives the following matrix modulo
$2^{20}=1048576$:
\[
 S\equiv
 \begin{pmatrix}
 917951&442211&476112&551248&968980\\
 679904&45903&922072&654544&644616\\
 797935&540216&1028224&572168&196164\\
 491016&385277&695495&798931&823817\\
 917289&206367&412992&727792&439172
 \end{pmatrix}\pmod {1048576}.
\]
The determinant of this displayed integer matrix is
$524288=2^{19}$ modulo $2^{20}$. Hence
\[
 v_2(\det S)=19,\qquad
 v_2(\det\mathcal A)=32,\qquad
 v_2(\det F)=26,\qquad v_2(\rho_{5,962})=6.
\]
The value $v_2(\det F)=26$ follows immediately from
\eqref{mt10:Fdet}: its only even factors beyond $2^{14}$ are
$(b-2)^2=960^2$, of valuation $12$.
An independent check evaluates $\det\mathcal A/\det F$ directly
as a rational number; it gives the same valuation. The original
integer numerator also satisfies
$P_{10}(962;2B_{1922},\ldots,2B_{1906})\equiv65536\pmod {131072}$.

This is a finite nonvanishing certificate, not a numerical
approximation. It closes the only missing index below the sufficient
threshold. Together with the preceding congruence and dominance
arguments, it proves Theorem~\ref{mt10:main}.

\subsection{Ordinary convergent identities supplied by the completed family}

For $b\ge6$ and $0\le k<b$, let
$\lambda^{(b,k)}_t=(M_{5,b}^{-1})_{kt}$, $0\le t<b$, and define
the finite odd polynomial and its centered difference
\[
 K_{b,k}(u)=\sum_{t=0}^{b-1}\lambda^{(b,k)}_tu^{2b+2t+1},
 \qquad Q_{b,k}(x)=K_{b,k}(x+\tfrac12)-K_{b,k}(x-\tfrac12).
\]
The centered difference is precisely the coefficient transform
\eqref{mt10:C}, so its continued moment sum equals $Z_k$.
For the ordinary formulation put
\[
 \eta_{s,0}=\frac1{s-1},\qquad
 \eta_{s,j}=\frac{(s)_{2j-1}B_{2j}(1/2)}{(2j)!}\quad(j\ge1),
 \qquad c_{b,\ell}=\sum_{j=0}^{\ell}\eta_{10,j}\eta_{2b,\ell-j}.
\]
Let $T_{b,k}(x)$ be the polynomial consisting of the nonnegative
powers in
\[
             Q_{b,k}(x)x^{-2b-8}\sum_{\ell\ge0}c_{b,\ell}x^{-2\ell}.
\]
Only finitely many coefficients enter it, and its degree is at most
$2b-10$. Then the completed family gives the exact identities
\begin{equation}\label{mt10:ordinary}
 \boxed{\displaystyle
 \sum_{n=0}^{\infty}\left\{
 Q_{b,k}(x_n)
  \frac{\psi^{(9)}(n+1)\psi^{(2b-1)}(n+1)}{9!\,(2b-1)!}
       -T_{b,k}(x_n)\right\}=\zeta(2k).}
\end{equation}
The summands are $O(x_n^{-2})$. Every removed nonnegative even
power has centered continued sum $\zeta(-2j,1/2)=0$, so the
ordinary absolutely convergent sum has exactly the same value as
the moment combination. Formula \eqref{mt10:ordinary} is a finite
rational algorithm at each $b,k$, with nonsingularity now proved
for the entire parameter range. At $k=0$ its value is the rational
number $-1/2$.

\paragraph{Status and further research.}
The new logical step is the exact $b=962$ certificate and the
combination of the two independently audited ranges. The source
congruence, the finite contraction, and the eventual dominance
mechanism are retained with attribution. The same strategy is
available for smaller order twelve, but its finite cases and modular
determinant have not been proved here. A uniform argument for all
unequal even orders, or even the full gap-six family, is still a
substantive open problem. Higher-minor positivity cannot simply be
inferred from the known pairwise ratio inequalities; the determinant
sign already changes in some families.

\paragraph{Reproducible artifacts.}
\texttt{completion\_certificate.py} proves the missing case from
the displayed five-square matrix using no external project files.
\texttt{verify\_eventual\_ten.py} rederives all $25$ entry norms and
the exact comparison \eqref{mt10:base}.
\texttt{verify\_p10\_dependencies.py}, its attributed algebra helper,
and \texttt{p10\_dependency\_certificate.json} replay the full inherited
polynomial identity, modular substitution, and all small base matrices.
The three output records distinguish full replay from the deliberately
optional quick mode. None uses floating-point evidence for a theorem.
